%%%%%%%%%%  scan idx 109  |  folio 96  %%%%%%%%%%
\SGAPrintPage{96}%
\sgaStableTarget{sga6.e01.prop.2-2}\textbf{Proposition 2.2.} \uline{Let} \(F\in\Ob\Dpx(\underline{C}_S)\) \uline{and} \(n\in\mathbb Z\)\uline{. Then}:

a) \uline{When} \(F\) \uline{is strictly} \(n\)\uline{-pseudo-coherent (relative to} \(\underline{C}_0\)\uline{) and} \(u:E\to F\) \uline{is a quasi-isomorphism, the set of} \(X\) \uline{over} \(S\) \uline{such that there exists a quasi-isomorphism} \(v:F'\longrightarrow E_X\) \uline{with} \(F'\) \uline{strictly} \(n\)\uline{-pseudo-coherent is a refinement of} \(S\)\uline{.}

b) \uline{The following conditions are equivalent}:
\begin{enumerate}[label=(\roman*)]
\item \uline{the set of} \(X\) \uline{over} \(S\) \uline{such that there exists a quasi-isomorphism} \(u:E\longrightarrow F_X\) \uline{of} \(\Cpx(\underline{C}_X)\)\uline{, with} \(E\) \uline{strictly} \(n\)\uline{-pseudo-coherent, is a refinement of} \(S\);
\item[\((\mathrm{i}')\)] \uline{the same condition as (i), but with} \(u\) \uline{an isomorphism of} \(\Dpx(\underline{C}_X)\);
\item \uline{the set of} \(X\) \uline{over} \(S\) \uline{such that there exists an} \(n\)\uline{-quasi-isomorphism (1.4.3)} \(u:E\longrightarrow F_X\) \uline{of} \(\Cpx(\underline{C}_X)\)\uline{, with} \(E\) \uline{strictly perfect, is a refinement of} \(S\);
\item[\((\mathrm{ii}')\)] \uline{the same condition as (ii), but with} \(u\) \uline{an} \(n\)\uline{-isomorphism of} \(\Dpx(\underline{C}_X)\).
\end{enumerate}

\uline{Proof.} a) For an integer \(i\) sufficiently large, there exists an \(i\)-quasi-isomorphism \(v_i:F'_i\longrightarrow E\), with \(F'_i\) strictly \(n\)-pseudo-coherent (take \(F'_i=0\)). Let us show that, for \(i>n\), there exists, after refining \(S\), an \((i-1)\)-quasi-isomorphism \(v_{i-1}:F'_{i-1}\longrightarrow E\), with \(F'_{i-1}\) strictly \(n\)-pseudo-coherent. By (1.4.4), it suffices to show that \(\Ho^{i-1}(M')\), where \(M'\) is the cone of \(v_i\), is of \(\underline{C}_0\)-finite type. Now, with \(M''\) denoting the cone of \(uv_i\), the hypothesis that \(u\) is a quasi-isomorphism and (1.4.2) imply that
\[
  \Ho^j(M')\qis \Ho^j(M'')\qquad\text{for every }j
\]
and \(\Ho^j(M')=0\) for \(j\geq i\).
%%%%%%%%%%  scan idx 110  |  folio 97  %%%%%%%%%%
\SGAPrintPage{97}%
The cone formula \(M''{}^j=F'_i{}^{\,j+1}\oplus F^j\) shows, on the other hand, that \(M''\) is strictly \((i-1)\)-pseudo-coherent. Thus, by (2.1.1), \(\Ho^{i-1}(M'')=\Ho^{i-1}(M')\) is of finite type, whence our assertion. By iteration, and after refining \(S\), we deduce an \(n\)-quasi-isomorphism \(v_n:F'_n\longrightarrow E\), where \(F'_n\) is strictly \(n\)-pseudo-coherent. But, by (1.4.5), we can extend \(v_n\) to a quasi-isomorphism \(v:F'\longrightarrow E\) without changing anything in degree \(\geq n\), hence in particular with \(F'\) strictly \(n\)-pseudo-coherent, which proves a).

b) The equivalence of conditions (i) and \((\mathrm{i}')\) follows at once from a) and the definition of the isomorphisms of \(\Dpx^-(\underline{C}_X)\). For \(E\in\Ob\Dpx^-(\underline{C}_S)\) strictly \(n\)-pseudo-coherent, there exists an exact sequence
\[
  0\longrightarrow E'\longrightarrow E\longrightarrow E''\longrightarrow 0,
\]
where \(E'\) is strictly perfect, and \(E''{}^i=0\) for \(i\geq n\); consequently \(E'\to E\) is an \(n\)-quasi-isomorphism, giving the implication \((\mathrm{i})\Rightarrow(\mathrm{ii})\). The implication \((\mathrm{ii})\Rightarrow(\mathrm{ii}')\) is trivial, so it remains to prove \((\mathrm{ii}')\Rightarrow(\mathrm{i})\). Let, then, \(u:E\longrightarrow F\) be an \(n\)-isomorphism of \(\Dpx^-(\underline{C}_S)\), with \(E\) strictly perfect; \(u\) is therefore defined by a diagram of \(\Cpx^-(\underline{C}_S)\)
\[
\begin{tikzcd}
& E' \arrow[dl,"s"'] \arrow[dr,"u'"]\\
E && F ,
\end{tikzcd}
\]
where \(s\) is a quasi-isomorphism and \(u'\) an \(n\)-quasi-isomorphism. By a), there exists, after refining \(S\), a quasi-isomorphism \(t:E''\longrightarrow E'\), with \(E''\) strictly \(n\)-pseudo-coherent. We then apply (1.4.5) to extend \(u't:E''\to F\) to a quasi-isomorphism \(E'''\longrightarrow F\) with \(E'''\) strictly \(n\)-pseudo-coherent, which proves \((\mathrm{ii}')\Rightarrow(\mathrm{i})\) and completes the proof of (2.2).

%%%%%%%%%%  scan idx 111  |  folio 98  %%%%%%%%%%
\SGAPrintPage{98}%
\sgaStableTarget{sga6.e01.def.2-3}\textbf{Definition 2.3.} \uline{Under the hypotheses of (2.2), we shall say that} \(F\in\Ob\Dpx(\underline{C}_S)\) \uline{is} \(\underline{n}\)-\(\underline{C}_0\)-pseudo-coherent \uline{when} \(F\) \uline{satisfies the equivalent conditions b) of (2.2). We shall say that} \(F\) \uline{is} \(\underline{C}_0\)-pseudo-coherent \uline{when} \(F\) \uline{is} \(n\)\uline{-}\(\underline{C}_0\)\uline{-pseudo-coherent for every} \(n\)\uline{.}

When this will cause no confusion, we shall sometimes allow ourselves to say \(n\)-pseudo-coherent, pseudo-coherent, instead of \(n\)-\(\underline{C}_0\)-pseudo-coherent, \(\underline{C}_0\)-pseudo-coherent.

\sgaStableTarget{sga6.e01.rem.2-3-1}\textbf{Remark 2.3.1.} Let \(F\in\Ob\Dpx(\underline{C}_S)\) be a pseudo-coherent complex. The set of \(X\) over \(S\) such that there exists a quasi-isomorphism \(E\longrightarrow F_X\) with \(E\) strictly pseudo-coherent (2.1) is in general not a refinement of \(S\). However, we shall see later (2.7. a) and Exp.II) important cases where it is.

The proof of (2.2 a)) gives the following result:

\sgaStableTarget{sga6.e01.prop.2-4}\textbf{Proposition 2.4.} \uline{Let} \(n,i\in\mathbb Z\)\uline{, and let} \(u:E\to F\) \uline{be an} \(i\)\uline{-quasi-isomorphism of} \(\Cpx(\underline{C}_S)\)\uline{, with} \(F\) \(n\)\uline{-pseudo-coherent and} \(E\) \uline{strictly} \(n\)\uline{-pseudo-coherent. Then, after refining} \(S\)\uline{, there exists a commutative triangle of} \(\Cpx(\underline{C}_S)\)
\[
\begin{tikzcd}
E \arrow[r,"v"] \arrow[d,"u"'] & E' \arrow[dl,"u'"]\\
F & ,
\end{tikzcd}
\]
\uline{where} \(E'\) \uline{is strictly} \(n\)\uline{-pseudo-coherent,} \(u'\) \uline{a quasi-isomorphism and} \(v\) \uline{a monomorphism inducing the identity in degree} \(\geq i\)\uline{.}

\sgaStableTarget{sga6.e01.prop.2-5}\textbf{Proposition 2.5.} a) \uline{When} \(F\in\Ob\Dpx(\underline{C}_S)\) \uline{is} \(n\)\uline{-pseudo-coherent (resp. pseudo-coherent),} \(F[m]\) \uline{is} \((n-m)\)\uline{-pseudo-coherent (resp. pseudo-coherent).}

%%%%%%%%%%  scan idx 112  |  folio 99  %%%%%%%%%%
\SGAPrintPage{99}%
b) \uline{Let}
\sgaStableTarget{sga6.e01.eq.unnamed.occ2}%
\[
\begin{tikzcd}
& G \arrow[dl,"+1"'] & \\
E \arrow[rr] && F \arrow[ul]
\end{tikzcd}
\tag{*}
\]
\uline{be a distinguished triangle of} \(\Dpx(\underline{C}_S)\)\uline{. When} \(E\) \uline{and} \(G\) \uline{are} \(n\)\uline{-pseudo-coherent (resp. pseudo-coherent), then} \(F\) \uline{is} \(n\)\uline{-pseudo-coherent (resp. pseudo-coherent).}

\uline{Proof.} a) Trivial from the definitions.

b) It of course suffices to prove the assertion without the ``resp.'' alternatives. The triangle \((*)\) can also be written
\[
\begin{tikzcd}
& F \arrow[dl,"+1"'] & \\
G[-1] \arrow[rr] && E \arrow[ul] & .
\end{tikzcd}
\]
By a), \(G[-1]\) is \((n+1)\)-pseudo-coherent. We are therefore reduced to proving that, when \(E\) is \((n+1)\)-pseudo-coherent and \(F\) \(n\)-pseudo-coherent, then \(G\) is \(n\)-pseudo-coherent. After localizing, we can assume that \(E\) (resp.\ \(F\)) is strictly \((n+1)\)-pseudo-coherent (resp.\ \(n\)-pseudo-coherent). The arrow \(E\xrightarrow{u}F\) of \(\Dpx^-(\underline{C}_S)\) is defined by a pair of arrows of \(\Cpx^-(\underline{C}_S)\)
\[
\begin{tikzcd}
& E' \arrow[dl,"s"'] \arrow[dr,"u'"]\\
E && F ,
\end{tikzcd}
\]
where \(s\) is a quasi-isomorphism. By (2.2 a)), there exists, after localizing, a quasi-isomorphism \(t:E''\to E'\), where \(E''\) is strictly \((n+1)\)-pseudo-coherent. We are therefore reduced to the case where \(u\) is defined by an
%%%%%%%%%%  scan idx 113  |  folio 100  %%%%%%%%%%
\SGAPrintPage{100}%
``actual'' arrow of \(\Cpx^-(\underline{C}_S)\). The assertion then follows from the definitions and the formula \(C(u)^i=E^{i+1}\oplus F^i\), where \(C(u)\) denotes the cone of \(u\).

\sgaStableTarget{sga6.e01.schol.2-6}\textbf{Scholium 2.6.} By ``rotation'' of the triangles, we see that pseudo-coherence properties of two of the objects of the triangle \((*)\) imply another such property of the third. The different cases can be summarized in the table below (the lowercase letters denote degrees of pseudo-coherence and each row corresponds to an implication whose conclusion is underlined):
\[
\begin{array}{ccc}
E&F&G\\[0.2em]
n+1&n&\underline n\\
n&\underline n&n\\
\underline n&n&n-1 .
\end{array}
\]

We shall denote by \(\Dpx(\underline{C})_{\underline{C}_0\text{-}n\text{-coh}}\) (resp.\ \(\Dpx(\underline{C})_{\underline{C}_0\text{-}\mathrm{coh}}\)) the strictly full fibered subcategory of \(\Dpx(\underline{C})\) whose fibre at each \(X\in\Ob \underline{S}\) consists of the objects of \(\Dpx(\underline{C}_X)\) that are \(n\)-pseudo-coherent (resp.\ pseudo-coherent) relative to \(\underline{C}_0\). By (2.5), the category \(\Dpx(\underline{C})_{\underline{C}_0\text{-}\mathrm{coh}}\) is a \uline{triangulated} sub-\(\underline{S}\)\uline{-category} of \(\Dpx(\underline{C})\) (1.1).

\sgaStableTarget{sga6.e01.prop.2-7}\textbf{Proposition 2.7.} \uline{Assume that} \(\underline{S}\) \uline{is a punctual site, so that} \(\underline{C}=\underline{C}_S\)\uline{,} \(\underline{C}_0=\underline{C}_{0S}\)\uline{. Then}:

a) \uline{When} \(F\in\Ob\Dpx(\underline{C})\) \uline{is pseudo-coherent, there exists a quasi-isomorphism} \(E\longrightarrow F\)\uline{, with} \(E\) \uline{strictly pseudo-coherent} (2.1).

b) \uline{Denote by} \(\Kpx^{-,\varnothing}(\underline{C}_0)\) \uline{the full triangulated subcategory of} \(\Kpx^-(\underline{C}_0)\) \uline{consisting of the acyclic complexes. The canonical functor}
\[
  \Kpx^-(\underline{C}_0)/\Kpx^{-,\varnothing}(\underline{C}_0)
  \longrightarrow \Dpx(\underline{C})
\]
